% 《物理实验报告模板.tex》的共同版式；所有实验模板均载入本文件。
\documentclass[UTF8,a4paper,10pt,fontset=fandol]{ctexart}
\usepackage[left=3.175cm,right=3.175cm,top=2.54cm,bottom=2.54cm]{geometry}
\usepackage{graphicx,tikz,pgfplots,needspace,adjustbox,booktabs,fancyhdr,lastpage,amsmath,amssymb,array,tabularx,longtable}
\pgfplotsset{compat=1.18}
\setlength{\parindent}{0pt}
\setlength{\parskip}{2pt}
\setlength{\headheight}{18pt}
\linespread{1}
\newcommand{\SetField}[2]{\expandafter\def\csname reportfield@#1\endcsname{#2}}
\makeatletter
\newcommand{\Field}[2][1.35cm]{%
  \ifcsname reportfield@#2\endcsname
    \expandafter\ifx\csname reportfield@#2\endcsname\@empty
      \underline{\makebox[#1]{\strut}}%
    \else\csname reportfield@#2\endcsname\fi
  \else\underline{\makebox[#1]{\strut}}\fi}
\makeatother
\newcommand{\MathField}[2][1.35cm]{%
  \ifcsname reportfield@math_#2\endcsname
    \csname reportfield@math_#2\endcsname
  \else\text{\Field[#1]{#2}}\fi}
\newcommand{\coverfield}[2]{%
  \noindent{\heiti\fontsize{18pt}{48pt}\selectfont
  \makebox[3.3cm][l]{#1}\underline{\makebox[8.8cm][c]{\Field[0pt]{#2}}}\par}}
\newcommand{\bindingline}[1][2.85]{%
  \begin{tikzpicture}[remember picture,overlay]
    \foreach \i in {0,...,38}{%
      \pgfmathsetmacro{\bindingy}{-#1-0.52*\i}%
      \ifnum\i=13\relax
        \node[inner sep=0pt] at ([xshift=1.05cm,yshift=\bindingy cm]current page.north west) {\songti\fontsize{10.5pt}{10.5pt}\selectfont 装};
      \else\ifnum\i=19\relax
        \node[inner sep=0pt] at ([xshift=1.05cm,yshift=\bindingy cm]current page.north west) {\songti\fontsize{10.5pt}{10.5pt}\selectfont 订};
      \else\ifnum\i=25\relax
        \node[inner sep=0pt] at ([xshift=1.05cm,yshift=\bindingy cm]current page.north west) {\songti\fontsize{10.5pt}{10.5pt}\selectfont 线};
      \else
        \pgfmathsetmacro{\bindingtop}{\bindingy+0.20}%
        \pgfmathsetmacro{\bindingbottom}{\bindingy-0.20}%
        \draw[line width=0.35pt,dash pattern=on 1.2pt off 1.8pt]
          ([xshift=1.05cm,yshift=\bindingtop cm]current page.north west)
          -- ([xshift=1.05cm,yshift=\bindingbottom cm]current page.north west);
      \fi\fi\fi}
  \end{tikzpicture}%
}
% 装订线印在每一页正文和原始数据页，续页也沿用母模板样式。
\newif\ifreportbinding
\AddToHook{shipout/foreground}{\ifreportbinding\bindingline\fi}
\fancypagestyle{report}{\fancyhf{}\fancyfoot[C]{第\thepage 页\quad 共\pageref{LastPage}页}\renewcommand{\headrulewidth}{0.4pt}}
\fancypagestyle{record}{\fancyhf{}\fancyhead[C]{原始数据记录单}\fancyfoot[C]{第\thepage 页\quad 共\pageref{LastPage}页}\renewcommand{\headrulewidth}{0.4pt}}
\newcommand{\ReportCover}[1]{%
  \thispagestyle{empty}\vspace*{1.0cm}
  \begin{center}
    \includegraphics[width=8.184cm]{tongji_logo.jpeg}\par
    {\bfseries\fontsize{18pt}{24pt}\selectfont TONGJI UNIVERSITY\par}
    \vspace{1.5cm}{\heiti\fontsize{36pt}{44pt}\selectfont 物理实验报告\par}
  \end{center}
  \vspace{1.5cm}
  \coverfield{实验名称}{experiment}
  \coverfield{学生姓名}{student}
  \coverfield{学\hspace{1em}号}{student_id}
  \coverfield{组\hspace{1em}号}{group}
  \coverfield{实验日期}{date}
  \clearpage\reportbindingtrue\pagestyle{report}
  \begin{center}{\heiti\fontsize{14pt}{40pt}\selectfont #1\par}\end{center}}
\newcommand{\ReportSection}[1]{\Needspace{4\baselineskip}\par\vspace{0.35\baselineskip}{\heiti #1}\par}
\newcommand{\ReportSubsection}[1]{\Needspace{3\baselineskip}\par\vspace{0.2\baselineskip}{\bfseries #1}\par}
\newif\ifreportimages
\newcommand{\ReportRecord}{%
  \clearpage\thispagestyle{record}
  \vspace*{0.6cm}%
  \ifreportimages\else（请将教师签字的实验数据记录单完整拍照粘贴至此处。）\fi
}
